\section{Methodology}
\label{sec:methodology}

\subsection{Problem Formulation}

We formulate bidirectional alignment as multi-objective RLHF:
\begin{equation}
\max_\theta \mathbb{E}_{x \sim \mathcal{D}, y \sim \pi_\theta(\cdot|x)} \left[ R_{\text{combined}}(x, y) - \beta_{\text{KL}} D_{\text{KL}}(\pi_\theta \| \pi_{\text{ref}}) \right]
\end{equation}
where:
\begin{equation}
R_{\text{combined}}(x, y) = \alpha \cdot R_{\text{help}}(x, y) + \beta \cdot R_{\text{ctrl}}(x, y)
\end{equation}

\subsection{IFEval as Differentiable Training Signal}

IFEval defines 25 verifiable constraint types. Binary constraint checks are non-differentiable; we convert to continuous scores via sigmoid soft thresholds:
\begin{equation}
s_i = \sigma\left(\frac{v_i - t_i}{\tau}\right)
\end{equation}

The controllability reward aggregates soft constraint scores:
\begin{equation}
R_{\text{ctrl}}(x, y) = \frac{1}{|C(x)|} \sum_{c_i \in C(x)} s_i
\end{equation}

We validated gradient flow in H-E1: variance = 0.039, scale.grad = 274.12.

\subsection{Experimental Conditions}

\paragraph{Treatments:} T1 ($\alpha$=0.2), T2 ($\alpha$=0.4), T3 ($\alpha$=0.6), T4 ($\alpha$=0.8)

\paragraph{Baselines:} B1 (SFT-only), B2 (Helpfulness RLHF), B3 (Quality RLHF)

\paragraph{Data Split:} IFEval 70/30 (training/held-out)
